\chapter{量子统计}

\section{Gibbs 因子}
在 6.1 节我们讨论系统与热库交换能量，现在我们允许交换粒子 $\Rightarrow dN\neq 0$。

则
\begin{align}
S_R(S_2)-S_R(S_1)
=-\frac{1}{T}\Bigl[E(S_2)-E(S_1)-\mu N(S_2)+\mu N(S_1)\Bigr].
\end{align}
于是
\begin{align}
\frac{P(S_2)}{P(S_1)}
=\frac{e^{-\left[E(S_2)-\mu N(S_2)\right]/kT}}{e^{-\left[E(S_1)-\mu N(S_1)\right]/kT}}
\quad\Rightarrow\quad
P(S)\propto e^{-\left[E(S)-\mu N(S)\right]/kT}
\qquad (\text{Gibbs 因子}).
\end{align}
归一化得
\begin{align}
P(S)=\frac{1}{\Xi}\,e^{-\left[E(S)-\mu N(S)\right]/kT},
\qquad
\Xi=\sum_S e^{-\left[E(S)-\mu N(S)\right]/kT}
\qquad (\text{巨配分函数}).
\end{align}

以 $\mathrm{O_2}$ 结合血红素为例：系统有两个态：未占据 $0$，与占据 $\varepsilon$。
\begin{align}
P(\mathrm{O_2})
=\frac{1}{\Xi}e^{-(\varepsilon-\mu)/kT}
=\frac{e^{-(\varepsilon-\mu)/kT}}{e^{-(\varepsilon-\mu)/kT}+1}.
\end{align}
又
\begin{align}
e^{-\mu/kT}=\frac{Z_{\mathrm{int}}}{Nv_Q}=\frac{kT\,Z_{\mathrm{int}}}{PV\,v_Q}.
\end{align}
因此
\begin{align}
P(\mathrm{O_2})
=\frac{1}{P_0/P+1}
=\frac{P}{P_0+P},
\qquad
P_0=\frac{kT\,Z_{\mathrm{int}}}{v_Q}e^{\varepsilon/kT}
\qquad (\text{Langmuir 吸附}).
\end{align}

\section{玻色子和费米子}
经典的统计
\begin{align}
Z=\frac{1}{N!}Z_1^N
\end{align}
的前提是粒子不可区分，即多个粒子占据相同单粒子态的可能性忽略不计（$Z_1\gg N$）。例如对于理想气体
\begin{align}
Z_1=\frac{V}{v_Q}Z_{\mathrm{int}},
\qquad
v_Q=\ell_Q^3=\left(\frac{h}{\sqrt{2\pi mkT}}\right)^3,
\end{align}
大约是平均 de Broglie 波长的三方。
\begin{align}
Z_1\gg N
\iff
\frac{V}{N}\gg v_Q,
\end{align}
即平均距离 $\gg$ de Broglie 波长，故是经典的。

但当 $Z_1\gg N$ 被违反时，多个粒子占据相同单粒子态不可忽略（变为量子系统）。此时粒子的性质很重要：可以与相同物质共享态的粒子为玻色子 (boson)，不能的叫做费米子 (fermion)。前者有整数自旋，后者有半整数自旋。

\subsection*{分布函数}
首先考虑由一个单粒子态组成的系统，这个态的能量为 $\varepsilon$，则它被 $n$ 个粒子占据的概率为
\begin{align}
P(n)=\frac{1}{\Xi}e^{-n(\varepsilon-\mu)/kT}.
\end{align}

对于费米子，$n=0/1$，
\begin{align}
\Xi=1+e^{-(\varepsilon-\mu)/kT},
\end{align}
故该态的占据数为
\begin{align}
\bar n_{\mathrm{FD}}
=\sum_n nP(n)=0\cdot P(0)+1\cdot P(1)
=\frac{e^{-(\varepsilon-\mu)/kT}}{1+e^{-(\varepsilon-\mu)/kT}}
=\frac{1}{e^{(\varepsilon-\mu)/kT}+1}
\qquad (\text{Fermi--Dirac 分布}).
\end{align}

对于玻色子，$n=0,1,2,\dots$，
\begin{align}
\Xi=\sum_{n=0}^{\infty}e^{-n(\varepsilon-\mu)/kT}
=\frac{1}{1-e^{-(\varepsilon-\mu)/kT}}.
\end{align}
并且
\begin{align}
\bar n_{\mathrm{BE}}
&=\sum_{n=0}^{\infty} nP(n)
=\sum_{n=0}^{\infty} n\cdot\frac{1}{\Xi}e^{-n(\varepsilon-\mu)/kT}
 =-\frac{1}{\Xi}\sum_{n=0}^{\infty}\frac{\partial}{\partial x}e^{-nx}
\quad\left(x=\frac{\varepsilon-\mu}{kT}\right)\\
&=-\frac{1}{\Xi}\frac{\partial \Xi}{\partial x}
=-(1-e^{-x})\frac{\partial}{\partial x}(1-e^{-x})^{-1}
=(e^{x}-1)^{-1}
=\frac{1}{e^{(\varepsilon-\mu)/kT}-1}
\qquad (\text{Bose--Einstein 分布}).
\end{align}

另外
\begin{align}
\bar n_{\mathrm{Boltzmann}}=NP(s)=\frac{N}{Z_1}e^{-\varepsilon/kT},
\qquad
\mu=-kT\ln\frac{Z_1}{N}\ (\text{参考习题 6.44}),
\end{align}
从而
\begin{align}
\bar n_{\mathrm{Boltzmann}}=e^{-(\varepsilon-\mu)/kT}.
\end{align}
当 $(\varepsilon-\mu)/kT\gg 1$ 时，$\bar n_{\mathrm{FD}}=\bar n_{\mathrm{BE}}=\bar n_{\mathrm{Boltzmann}}$。

\section{简并 Fermi 气体}
考虑极低温下的 Fermi 气体（如电子气体），这里极低温指 $V/N\ll v_Q$。

\subsection{温度为零}
$T=0$ 时，Fermi--Dirac 分布退化为阶跃函数，所有能量小于 $\mu$ 的单粒子态已被占据，其余均未被占据。此时 Fermi 气体是简并的，$\mu$ 被称为 Fermi 能，$\varepsilon_F\equiv \mu$。

$\varepsilon_F$ 代表电子的最高能量。为计算 $\varepsilon_F$，假设电子限制在 $V=L^3$ 内的自由粒子，则
\begin{align}
\lambda_n=\frac{2L}{n},\qquad
p_n=\frac{h}{\lambda_n}=\frac{hn}{2L}
\quad\Rightarrow\quad
\varepsilon=\frac{|\vec p|^2}{2m}=\frac{h^2}{8mL^2}(n_x^2+n_y^2+n_z^2).
\end{align}
可以想象一个能量空间，每个格点代表一个能量的两个态（自旋）。$N$ 足够大时，所有电子占据的格点近似构成一个球面。

记半径为 $n_{\max}$，则
\begin{align}
\varepsilon_F=\frac{h^2n_{\max}^2}{8mL^2}.
\end{align}
另一方面
\begin{align}
N=2\times\frac{1}{8}\times\frac{4}{3}\pi n_{\max}^3=\frac{\pi n_{\max}^3}{3}
\quad\Rightarrow\quad
n_{\max}=\left(\frac{3N}{\pi}\right)^{1/3}.
\end{align}
故
\begin{align}
\varepsilon_F
=\frac{h^2}{8mL^2}\left(\frac{3N}{\pi}\right)^{2/3}
=\frac{h^2}{8m}\left(\frac{3N}{\pi V}\right)^{2/3}.
\end{align}
这是一个强度量，只与电子密度 $N/V$ 有关 $\Rightarrow$ 适用于任何宏观气态系统。

电子的平均能量会小一些：
\begin{align}
U
&=2\sum_{n_x}\sum_{n_y}\sum_{n_z}\varepsilon(\vec n)
=2\iiint \varepsilon(\vec n)\,dn_x\,dn_y\,dn_z
=2\int_0^{n_{\max}}dn\int_0^{\pi/2}d\theta\int_0^{\pi/2}d\phi\,n^2\sin\theta\,\varepsilon(n)\\
&=\pi\int_0^{n_{\max}}\varepsilon(n)n^2\,dn
=\pi\frac{h^2}{8mL^2}\int_0^{n_{\max}}n^4\,dn
=\frac{\pi h^2n_{\max}^5}{40mL^2}
=\frac{3}{5}N\varepsilon_F
\quad\Rightarrow\quad
\bar\varepsilon=\frac{3}{5}\varepsilon_F.
\end{align}

对于典型金属中的电子，$\varepsilon_F\sim 3\,\mathrm{eV}$，而室温下 $kT\sim \tfrac{1}{40}\,\mathrm{eV}$，故 $V/N\ll v_Q\iff kT\ll \varepsilon_F$，此时作 $T\approx 0$ 近似是合理的（涨落可忽略）。

对于 Fermi 气体，使 $kT=\varepsilon_F$ 的温度为 Fermi 温度
\begin{align}
T_F=\frac{\varepsilon_F}{k}.
\end{align}
同时可以计算 Fermi 气体的压强（简并压），这个压强来自不相容原理：
\begin{align}
P=-\frac{\partial U}{\partial V}
=-\frac{\partial}{\partial V}\left(\frac{3}{5}N\varepsilon_F\right)
=\frac{2N\varepsilon_F}{5V}=\frac{2}{3}\frac{U}{V}.
\end{align}
这个值 $P>0$，对抗金属中的静电力。

实际可测量的是压缩模量：
\begin{align}
B=-V\left(\frac{\partial P}{\partial V}\right)_T=\frac{10}{9}\frac{U}{V}.
\end{align}

\subsection{较小的 $T\neq 0$}
在 $T$ 下，所有粒子通常会获得 $kT$ 热能，但在简并 Fermi 气体中，大部分粒子获得 $kT$ 而可能进入的态已被占据，能够获得 $kT$ 的粒子只有处于靠近 $\varepsilon_F$ 的于 $kT$ 的粒子。
\begin{align}
\Rightarrow\ \text{受 $T$ 增加影响的电子数}\propto T,\qquad \text{同时也}\propto N\ (\text{广延量}).
\end{align}
因此额外的能量
\begin{align}
\Delta U\propto (NkT)(kT)=N(kT)^2.
\end{align}
从量纲上看，$[N(kT)^2]=\mathrm{J}^2$，需除以一个单位为能量的量 $\Rightarrow\varepsilon_F$。最后得到（精确证明）
\begin{align}
U=\frac{3}{5}N\varepsilon_F+\frac{\pi^2}{4}N\frac{(kT)^2}{\varepsilon_F}
\quad\Rightarrow\quad
C_V=\left(\frac{\partial U}{\partial T}\right)_V=\frac{\pi^2Nk^2T}{2\varepsilon_F}.
\end{align}

\subsection{态密度}
对
\begin{align}
U=\pi\int_0^{n_{\max}}\varepsilon(n)n^2\,dn
\end{align}
作换元：
\begin{align}
\varepsilon=\frac{h^2}{8mL^2}n^2
\quad\Rightarrow\quad
n=\sqrt{\frac{8mL^2}{h^2}}\,\sqrt{\varepsilon},
\qquad
dn=\sqrt{\frac{8mL^2}{h^2}}\,\frac{1}{2\sqrt{\varepsilon}}\,d\varepsilon.
\end{align}
得到 $T=0$ 时
\begin{align}
U=\int_0^{\varepsilon_F}\varepsilon\left[\frac{\pi}{2}\left(\frac{8mL^2}{h^2}\right)^{3/2}\right]\sqrt{\varepsilon}\,d\varepsilon.
\end{align}
括号内的记作态密度 $g(\varepsilon)$：
\begin{align}
g(\varepsilon)
&=\frac{\pi}{2}\left(\frac{8mL^2}{h^2}\right)^{3/2}\sqrt{\varepsilon}
=\frac{\pi(8m)^{3/2}}{2h^3}V\sqrt{\varepsilon}
=\frac{3N}{2\varepsilon_F^{3/2}}\sqrt{\varepsilon}
=g_0\sqrt{\varepsilon}\propto \sqrt{\varepsilon}.
\end{align}

对于 $T=0$ 的气体，
\begin{align}
N=\int_0^{\varepsilon_F}g(\varepsilon)\,d\varepsilon.
\end{align}
$T\neq 0$ 时则引入 Fermi--Dirac 分布：
\begin{align}
N&=\int_0^{\infty}g(\varepsilon)\,\bar n_{\mathrm{FD}}(\varepsilon)\,d\varepsilon
=\int_0^{\infty}g(\varepsilon)\,\frac{1}{e^{(\varepsilon-\mu)/kT}+1}\,d\varepsilon,\\
U&=\int_0^{\infty}\varepsilon g(\varepsilon)\,\bar n_{\mathrm{FD}}(\varepsilon)\,d\varepsilon
=\int_0^{\infty}\varepsilon g(\varepsilon)\,\frac{1}{e^{(\varepsilon-\mu)/kT}+1}\,d\varepsilon.
\end{align}
此时化学势（一个态被占据的概率为 1 的 $\varepsilon$）$\mu<\varepsilon_F$（参考习题 7.28）。

\subsection{Sommerfeld 展开}
我们希望计算 $N$，从而求解 $\mu(T)$，代入 $U(T)$，最终得到 $C_V(T)$，但它们无法直接积分求得。Sommerfeld 展开可以近似地计算：
\begin{align}
N
&=\int_0^{\infty}g(\varepsilon)\,\bar n_{\mathrm{FD}}(\varepsilon)\,d\varepsilon
=g_0\int_0^{\infty}\sqrt{\varepsilon}\,\bar n_{\mathrm{FD}}(\varepsilon)\,d\varepsilon
=\frac{2}{3}g_0\varepsilon^{3/2}\bar n_{\mathrm{FD}}(\varepsilon)\Big|_0^{\infty}+\frac{2}{3}g_0\int_0^{\infty}\varepsilon^{3/2}\left(-\frac{d\bar n_{\mathrm{FD}}}{d\varepsilon}\right)d\varepsilon\\
&=\frac{2}{3}g_0\int_0^{\infty}\frac{1}{kT}\frac{e^x}{(e^x+1)^2}\,\varepsilon^{3/2}\,d\varepsilon
=\frac{2}{3}g_0\int_{-\mu/kT}^{\infty}\frac{e^x}{(e^x+1)^2}\,\varepsilon^{3/2}\,dx
\qquad\left(x=\frac{\varepsilon-\mu}{kT}\right).
\end{align}
当 $|\varepsilon-\mu|\gg kT$ 时，$-\dfrac{d\bar n_{\mathrm{FD}}}{d\varepsilon}=\dfrac{1}{kT}\dfrac{e^x}{(e^x+1)^2}$ 指数衰减，故可作两处近似：把积分下限延展到 $-\infty$，并在 $\varepsilon=\mu$ 处 Taylor 展开 $\varepsilon^{3/2}$：
\begin{align}
\varepsilon^{3/2}=\mu^{3/2}+\frac{3}{2}(\varepsilon-\mu)\mu^{1/2}+\frac{3}{8}(\varepsilon-\mu)^2\mu^{-1/2}+\cdots.
\end{align}
于是
\begin{align}
N
&\approx \frac{2}{3}g_0\int_{-\infty}^{\infty}\frac{e^x}{(e^x+1)^2}\left[\mu^{3/2}+\frac{3}{2}xkT\mu^{1/2}+\frac{3}{8}(xkT)^2\mu^{-1/2}+\cdots\right]dx\\
&=\frac{2}{3}g_0\mu^{3/2}+\frac{\pi^2}{12}g_0(kT)^2\mu^{-1/2}+\cdots
= N\left(\frac{\mu}{\varepsilon_F}\right)^{3/2}+N\frac{\pi^2}{8}\frac{(kT)^2}{\varepsilon_F^{3/2}\mu^{1/2}}+\cdots.
\end{align}
消去 $N$，可以看出 $\mu/\varepsilon_F\approx 1$，修正项 $\propto (kT/\varepsilon_F)^2$ 很小，可以用 $\mu$ 替换 $\varepsilon_F$，解出
\begin{align}
\frac{\mu}{\varepsilon_F}=\left[1-\frac{\pi^2}{8}\left(\frac{kT}{\varepsilon_F}\right)^2+\cdots\right]^{2/3}
=1-\frac{\pi^2}{12}\left(\frac{kT}{\varepsilon_F}\right)^2+\cdots.
\end{align}
同样方法可以计算
\begin{align}
U=\frac{3}{5}N\varepsilon_F+\frac{\pi^2}{4}N\frac{(kT)^2}{\varepsilon_F}+\cdots.
\end{align}

\section{黑体辐射}
考虑一个量子化的电磁场，可以看作是不同波长的驻波叠加。每个驻波对应一个谐振子，其允许的能级为 $E_n=0,hf,2hf,\dots$，因此单个谐振子的配分函数
\begin{align}
Z_1=1+e^{-\beta hf}+e^{-2\beta hf}+\cdots=\frac{1}{1-e^{-\beta hf}}
\quad\Rightarrow\quad
\bar E=-\frac{1}{Z_1}\frac{\partial Z_1}{\partial \beta}=\frac{hf}{e^{hf/kT}-1}.
\end{align}
把 $hf$ 看作能量单位，就得到能量除数的平均值
\begin{align}
\bar n_{\mathrm{pl}}=\frac{1}{e^{hf/kT}-1}
\qquad (\text{Planck 分布}).
\end{align}

电磁场中的能量单位是光子（一种玻色子），故也满足 Bose--Einstein 分布。比较 $\bar n_{\mathrm{BE}}$ 和 $\bar n_{\mathrm{pl}}$，$\varepsilon=hf\Rightarrow\mu=0$。

这一点可以从光子自由产生/湮灭看出来：当 $T,V$ 固定，系统的 $F$ 一定取到最小。对于光子系统，$N$ 并不守恒，但 $F$ 始终最小，故
\begin{align}
\mu=\left(\frac{\partial F}{\partial N}\right)_{T,V}=0.
\end{align}

类似电子气体，现在对光子的模式求和。注意光子有两个偏振方向，即每个波矢对应 $2$ 个。

一维长度为 $L$ 的盒子中，
\begin{align}
\lambda=\frac{2L}{n}\Rightarrow p=\frac{h}{\lambda}=\frac{hn}{2L}\Rightarrow \varepsilon=pc=\frac{hcn}{2L}.
\end{align}
换到三维，
\begin{align}
\varepsilon=c|\vec p|=\frac{hc}{2L}\sqrt{n_x^2+n_y^2+n_z^2}=\frac{hcn}{2L}.
\end{align}

先求总数 $N$：
\begin{align}
N
&=2\sum_{n_x}\sum_{n_y}\sum_{n_z}\bar n_{\mathrm{pl}}(\varepsilon)
=2\sum_{n_x,n_y,n_z}\frac{1}{e^{hc n/(2LkT)}-1}
=\int_0^{\infty}dn\int_0^{\pi/2}d\theta\int_0^{\pi/2}d\phi\,\frac{n^2\sin\theta}{e^{hc n/(2LkT)}-1}\\
&=\pi\int_0^{\infty}\frac{n^2}{e^{hc n/(2LkT)}-1}\,dn
=\pi\left(\frac{2LkT}{hc}\right)^3\int_0^{\infty}\frac{x^2}{e^x-1}\,dx
\qquad\left(x=\frac{hcn}{2LkT}\right).
\end{align}
这个积分无法解析求解。

再求总能量：
\begin{align}
U
&=2\sum_{n_x}\sum_{n_y}\sum_{n_z}\varepsilon\,\bar n_{\mathrm{pl}}(\varepsilon)
=\sum_{n_x,n_y,n_z}\frac{hcn}{2L}\frac{1}{e^{hcn/(2LkT)}-1}\\
&=\frac{\pi}{2}\int_0^{\infty}\frac{hc}{L}\frac{n^3}{e^{hcn/(2LkT)}-1}\,dn
=V\int_0^{\infty}\frac{8\pi}{(hc)^3}\frac{\varepsilon^3}{e^{\varepsilon/kT}-1}\,d\varepsilon
\qquad\left(\varepsilon=\frac{hcn}{2L},\ L^3=V\right).
\end{align}
因此
\begin{align}
\frac{U}{V}=\int_0^{\infty}\frac{8\pi}{(hc)^3}\frac{\varepsilon^3}{e^{\varepsilon/kT}-1}\,d\varepsilon.
\end{align}
此时被积函数数量单位光子能量的能量密度，光谱为
\begin{align}
u(\varepsilon)=\frac{8\pi}{(hc)^3}\frac{\varepsilon^3}{e^{\varepsilon/kT}-1}.
\end{align}

再作一次换元得到无量纲变量 $x=\varepsilon/kT$：
\begin{align}
\frac{U}{V}
&=\frac{8\pi(kT)^4}{(hc)^3}\int_0^{\infty}\frac{x^3}{e^x-1}\,dx
=\frac{8\pi(kT)^4}{(hc)^3}\frac{\pi^4}{15}
=\frac{8\pi^5(kT)^4}{15(hc)^3}
\propto T^4.
\end{align}
进一步还可以计算其他量：
\begin{align}
C_V=\left(\frac{\partial U}{\partial T}\right)_V=4aT^3,
\qquad
a=\frac{8\pi^5k^4}{15(hc)^3}V.
\end{align}
\begin{align}
S(T)=\int_0^T\frac{C_V(T')}{T'}\,dT'
=4a\int_0^T {T'}^2\,dT'
=\frac{4}{3}aT^3
=\frac{32\pi^5}{45}V\left(\frac{kT}{hc}\right)^3k.
\end{align}
\begin{align}
P=-\left(\frac{\partial U}{\partial V}\right)_{S,N},\quad \text{为固定 }S,\ U=aVT^4=\left(\frac{3S}{4}\right)^{4/3}(aV)^{-1/3}
\end{align}
因此
\begin{align}
P=-\left(\frac{\partial U}{\partial V}\right)_{S,N}=\frac{U}{3V}=\frac{1}{3}aT^4.
\end{align}
并且
\begin{align}
F=U-TS=aVT^4-\frac{4}{3}aVT^4=-\frac{1}{3}aVT^4=-\frac{1}{3}U.
\end{align}

现在我们分析了处热平衡的室内光的气体，但实际上我们试图研究的是热物体发射的光子。为此，我们考虑围绕孔的小黑体试亮。

逃逸出去的能量 = 试亮总能量 $\times$ 逃逸概率：
\begin{align}
&=\int_0^{2\pi}d\phi\int_0^{\pi/2}d\theta\,\frac{U}{V}\,\bigl(Rd\theta\bigr)\bigl(R\sin\theta\,d\phi\bigr)\,c\,dt\times \frac{A\cos\theta}{4\pi R^2}\\
&=\int_0^{2\pi}d\phi\int_0^{\pi/2}\sin\theta\cos\theta\,d\theta\,\frac{AUc}{4\pi V}\,dt
=\frac{AU}{4V}c\,dt.
\end{align}
因此单位面积辐射功率
\begin{align}
\frac{1}{A}\frac{dU}{dt}=\frac{c}{4}\frac{U}{V}
=\frac{2\pi^5}{15}\frac{(kT)^4}{h^3c^2}=\sigma T^4.
\end{align}
Stefan--Boltzmann 常量
\begin{align}
\sigma=\frac{2\pi^5k^4}{15h^3c^2}=5.67\times 10^{-8}\ \mathrm{W\,m^{-2}\,K^{-4}}.
\end{align}

Stefan 定律也适用于温度为 $T$ 时由任何非反射表面辐射的光子（黑体辐射）。但若物体会反射一些光子，而对吸收剩下的光子比例为 $e$，则与完美黑体相似，$e$ 也是发射光子的比例。$e$ 称为发射率，通常取决于光的波长。若对所有波长使用 $e$ 的加权平均，则物体辐射总功率为
\begin{align}
\sigma eAT^4,
\end{align}
$A$ 为表面积。

关于太阳与地球的例子，参考 7.4.11 和习题 7.55。

\section{固体的 Debye 模型}
2.2 节讨论了固体的 Einstein 模型，这个模型将原子视为独立的三维谐振子。

Einstein 模型预测
\begin{align}
C_V=3Nk\frac{(\varepsilon/kT)^2e^{\varepsilon/kT}}{(e^{\varepsilon/kT}-1)^2}
\qquad (\text{习题 3.25}),
\end{align}
$N$ 是原子数，$\varepsilon=hf$ 是能量单位。

$kT\gg \varepsilon$ 时，$C_V\to 3Nk$ 符合能量均分定理；$T\to 0$ 时，模型预测 $C_V$ 指数衰减到 $0$，但实验测得 $C_V\propto T^3$。

这是因为晶体中的原子不是独立振动，在低温仍有许多低频振动模式未被冻结。它们类似于真空中的电磁波动振荡。这种在晶体中的机械振动称为声波，由声子介导。类似光子，我们也可以写出声子的 Planck 分布（$c_s$ 声速，$L$ 晶体长度）：
\begin{align}
\varepsilon=hf=\frac{hc_s}{\lambda}=\frac{hc_sn}{2L}\quad (n=|\vec n|\text{ 是声波矢量}),
\qquad
\bar n_{\mathrm{pl}}=\frac{1}{e^{\varepsilon/kT}-1}.
\end{align}

类似地，我们对所有模式求和，但这里求和上限不是 $\infty$，因为声波波长不能小于两个原子间距 $\Rightarrow n\le 3\sqrt{N}$。这里我们假设晶体是完美正方体，但这样空间并不连续积分（这样 $n_x,n_y,n_z$ 的关系很复杂）。一个好的近似是把 $\vec n$ 空间缩成体积为 $N$ 的高球体：
\begin{align}
n_{\max}=\left(\frac{6N}{\pi}\right)^{1/3}.
\end{align}

Debye 近似在高低温极限都很准确，高温时关键是求出总自由度，故体积一致即可；低温时 $|\vec n|$ 较大的模式被冻结，因此可以对它们近似。

总之
\begin{align}
U
&=3\sum_{n_x}\sum_{n_y}\sum_{n_z}\varepsilon\,\bar n_{\mathrm{pl}}(\varepsilon)
\qquad (3\text{ 来自每个 }\vec n\text{ 对应三个极化态})\\
&=3\int_0^{n_{\max}}dn\int_0^{\pi/2}d\theta\int_0^{\pi/2}d\phi\,\frac{n^2\sin\theta\,\varepsilon}{e^{\varepsilon/kT}-1}
=3\pi\int_0^{n_{\max}}\frac{hc_s}{2L}\frac{n^3}{e^{hc_sn/(2LkT)}-1}\,dn.
\end{align}
令
\begin{align}
x=\frac{hc_sn}{2LkT},
\qquad
x_{\max}=\frac{hc_sn_{\max}}{2LkT}
=\frac{hc_s}{2LkT}\left(\frac{6N}{\pi}\right)^{1/3}\equiv \frac{T_D}{T},
\end{align}
得到
\begin{align}
U=\frac{9NkT^4}{T_D^3}\int_0^{T_D/T}\frac{x^3}{e^x-1}\,dx.
\end{align}
这里引入 Debye 温度
\begin{align}
T_D=\frac{hc_s}{2k}\left(\frac{6N}{\pi V}\right)^{1/3}.
\end{align}
这个积分无法解析求解，但可以考虑高低温极限。

$T\gg T_D$ 时，$T_D/T\ll 1$，此时 $e^x-1\approx x$：
\begin{align}
U\approx \frac{9NkT^4}{T_D^3}\int_0^{T_D/T}x^2\,dx=3NkT
\quad\Rightarrow\quad
C_V\approx 3Nk.
\end{align}
$T\ll T_D$ 时，$T_D/T\to\infty$：
\begin{align}
U\approx \frac{9NkT^4}{T_D^3}\int_0^{\infty}\frac{x^3}{e^x-1}\,dx
=\frac{3\pi^4}{5}Nk\frac{T^4}{T_D^3}
\quad\Rightarrow\quad
C_V=\frac{12\pi^4Nk}{5T_D^3}T^3\propto T^3.
\end{align}
与实验有较好吻合。对于中间温度，直接对 $U=\dfrac{9NkT^4}{T_D^3}\int_0^{T_D/T}\dfrac{x^3}{e^x-1}\,dx$ 求导，得到
\begin{align}
C_V=9Nk\left(\frac{T}{T_D}\right)^3\int_0^{T_D/T}\frac{x^4e^x}{(e^x-1)^2}\,dx.
\end{align}
$T=T_D$ 时，$C_V$ 达到最大值的 $95\%$，这说明 $T_D$ 实际上代表一个我们可以大致能量均分定理的温度。

\section{Bose--Einstein 凝聚}
7.4、7.5 节我们处理了可以任意数量存在的玻色子（光子、声子），其总数由热平衡确定。现在我们想知道数量固定的普通玻色子是怎样的。

我们先来确定化学势，考虑 $T>0$ 的极限。在 $T=0$ 时，所有原子处在基态。对于处在 $V=L^3$ 中的原子，
\begin{align}
\varepsilon_0=\frac{h^2}{8mL^2}(1^2+1^2+1^2)=\frac{3h^2}{8mL^2}.
\end{align}
这个值相当小，意味着
\begin{align}
e^{(\varepsilon_0-\mu)/kT}-1\approx \frac{\varepsilon_0-\mu}{kT}.
\end{align}
因此在 $T$ 足够低时，基态上的平均原子数
\begin{align}
N_0=\frac{1}{e^{(\varepsilon_0-\mu)/kT}-1}=\frac{kT}{\varepsilon_0-\mu}
\qquad (N_0\gg 1\text{ 时}).
\end{align}
$\Rightarrow$ $T=0$ 时，$\mu=\varepsilon_0$；$T\neq 0$ 但很小时，$\mu\approx\varepsilon_0$ 且 $\mu<\varepsilon_0$。

那么 $T$ 要多低才能使 $N_0\gg 1$？考虑
\begin{align}
N=\sum_s\frac{1}{e^{(\varepsilon_s-\mu)/kT}-1}\approx \int_0^{\infty}g(\varepsilon)\frac{1}{e^{(\varepsilon-\mu)/kT}-1}\,d\varepsilon.
\end{align}
这里态密度是 7.3 节的一半（只有一种自旋）：
\begin{align}
g(\varepsilon)=\frac{1}{4\pi^2}\left(\frac{2m}{\hbar^2}\right)^{3/2}V\sqrt{\varepsilon}.
\end{align}
这个积分无法解析求解，因此为了确定 $\mu$，必须不断做数值积分。考虑 $\mu=0$，这是 $T$ 足够低、$N_0$ 足够大时一个好近似。代入 $\mu=0$，令 $x=\varepsilon/kT$：
\begin{align}
N
&=\frac{1}{4\pi^2}\left(\frac{2m}{\hbar^2}\right)^{3/2}V\int_0^{\infty}\frac{\sqrt{\varepsilon}\,d\varepsilon}{e^{\varepsilon/kT}-1}
=\frac{1}{4\pi^2}\left(\frac{2mkT}{\hbar^2}\right)^{3/2}V\int_0^{\infty}\frac{x^{1/2}}{e^x-1}\,dx\\
&=\zeta\!\left(\frac{3}{2}\right)\left(\frac{2\pi mkT}{h^2}\right)^{3/2}V
\approx 2.612\left(\frac{2\pi mkT}{h^2}\right)^{3/2}V.
\end{align}
反解出对应的温度
\begin{align}
kT_c=0.527\left(\frac{h^2}{2\pi m}\right)\left(\frac{N}{V}\right)^{2/3}.
\end{align}

$T\neq T_c$ 时会发生什么？$T>T_c$ 时，$\mu$ 一定是小于 $0$；$T<T_c$ 时，如果 $\mu\to 0$，$N$ 一定变得更大，但我们已经固定了 $N$，发生了什么？

问题在于把求和近似为积分的一步。在 $T\to 0$ 时，基态上有很多粒子，不能把离散的能级近似为连续的谱。实际上，应该把 $\varepsilon=0$ 的态单独看作离散基态，而把其余作为从第一个行态 $\varepsilon\sim kT$ 之间的下限开始，描述 $\varepsilon\gg \varepsilon_0$ 的大量态。

此时
\begin{align}
N_{\mathrm{excited}}=2.612\left(\frac{2\pi mkT}{h^2}\right)^{3/2}V
\qquad (T<T_c).
\end{align}
总之，$T<T_c$ 时，$\mu\approx 0$，
\begin{align}
N_{\mathrm{excited}}=\left(\frac{T}{T_c}\right)^{3/2}N
\quad\Rightarrow\quad
N_0=N-N_{\mathrm{excited}}=\left[1-\left(\frac{T}{T_c}\right)^{3/2}\right]N.
\end{align}

这种在 $T<T_c$ 时发生的基态原子宏观聚集称为 Bose--Einstein 凝聚，$T_c$ 为凝聚温度。回顾
\begin{align}
v_Q=\left(\frac{h}{\sqrt{2\pi mkT}}\right)^3,
\end{align}
$T=T_c$ 时 $v_Q=V/N$。也就是说当原子的热波动长度与平均间距相当时，凝聚就会发生。